\clearpage\begingroup\SourceLeaf{53b}\HistoricalHead{48}{L}{PHYSIK. BRIEFWECHSEL.}
\begin{center}
[2.]\\[.7em]
\textsc{Gauss} an \textsc{Olbers}. Göttingen, 14. Juni 1830.
\par\smallskip\rule{.72\textwidth}{.35pt}\\
{\small [\textsc{Wilhelm Olbers}, Sein Leben und seine Werke II 2, 1909, S. 544.]}\\
\rule{.72\textwidth}{.35pt}
\end{center}
\noindent .\ .\ .\ .\ .\ .\ .\ .\ .\ .\ .\ .\ .\ .\ .\ .\ .\ .\ .\ .\ .
\par\noindent\hspace*{\parindent}Von dem \textsc{Herzog von Sussex} hat die Sternwarte vor einiger Zeit ein von
\textsc{Hardy} verfertigtes verkehrtes Pendel erhalten, in einer etwa 4 mal so grossen
Dimension wie das früher von demselben geschenkte (G.G.A., 1820, S. 1866[\textsuperscript{*})]\begingroup\renewcommand{\thefootnote}{}\footnotetext{[\textsuperscript{*}) Werke VI, S. 435.]}\endgroup).
Es hat eine bedeutende Empfindlichkeit, wenn gleich \textsc{Hardys} Hoffnung, dass
man damit die Attraction der Sonne, des Mondes und der Berge soll messen
oder augenfällig machen können, theils viel zu sanguinisch, theils ohne Sinn
ist. Bei einer mittleren Stellung der Kugel entspricht ein Ausschlag des
Zeigers von 1 Linie einer Neigung von 24 Secunden, oder die Empfindlichkeit
ist so gross, wie die eines 34 mal so langen gewöhnlichen Pendels. Was mich
dabei am meisten interessirt, ist die Schärfe, mit der man die Dauer der
Vibrationen bestimmen kann, und die Regelmässigkeit, mit welcher die Ungleichheit
derselben der Temperatur folgt. Ich werde meine Versuche damit
vervielfältigen. Hier einige, wobei das Th[ermometer] sorgfältig notirt ist:
\begin{center}
\begin{tabular}{c|c}
\hline\hline
Therm. & Dauer einer\\
Centes. & Vibration\\
\hline\hline
\(15^0{,}8\) & \(3^{\mathrm{s}}{,}352\)\\
\(16\phantom{^0}{,}8\) & \(3\phantom{^{\mathrm{s}}}{,}364\)\\
\(18\phantom{^0}{,}2\) & \(3\phantom{^{\mathrm{s}}}{,}378\)\\
\(19\phantom{^0}{,}4\) & \(3\phantom{^{\mathrm{s}}}{,}395\)
\end{tabular}
\end{center}
\clearpage\SourceLeaf{54}\HistoricalHead{49}{R}{ENCKE AN GAUSS.}
\noindent\hspace*{\parindent}Natürlich ist die Temperaturänderung nicht, insofern sie die Länge der
Stange oder die Dichtigkeit der Luft afficirt (dieser Einfluss ist so gut wie
unmerklich), sondern insofern sie die Elasticität der Feder verändert, das
Agens, und ich glaube, dass man diesen, so viel ich weiss, bisher noch nie
quantitativ untersuchten Umstand zur Darstellung einer ganz neuen Art äusserst
empfindlicher Thermometer müsse benutzen können.\quad .\ .\ .\ .\ .\ .\ .\ .\ .\ .\ .\ .\ .\ .
\clearpage\endgroup
